\chapter{Linear coupling}
\label{ch:linear-coupling}

\section{Scope and conventions}

We restrict the discussion to the transverse canonical phase-space vector
\begin{equation}
 \mat z=(x,p_x,y,p_y)^T,
\end{equation}
with symplectic matrix
\begin{equation}
 \mat S=\diag(\mat J,\mat J),
 \qquad
 \mat J=\begin{pmatrix}0&1\\-1&0\end{pmatrix}.
\end{equation}
The one-turn map is assumed stable and symplectic.  Mode 1 is labeled horizontal-like and mode 2 vertical-like.  A superscript $*$ denotes complex conjugation.  Following the notation of the Xsuite physics manual, for a real $2\times2$ matrix $A$ we write
\begin{equation}
 \bar A\equiv-\mat J A^T\mat J
\end{equation}
for its symplectic conjugate; the overbar must not be confused with complex
conjugation.

The exact matrix relations below apply to a purely transverse $4\times4$ symplectic problem.  Longitudinal coupling, radiation damping, non-symplectic maps, and the distinction between canonical momentum and geometric slope require additional terms.  The complex-variable and RDT signs follow the convention introduced in \cref{ch:rdt}.

\section{Turn-by-turn motion in the Mais--Ripken formalism}
\label{sec:mr}

\subsection{Linear normal form and the $W$ matrix}

At observation point $s$, let $R(s)$ be the one-turn map.  Its real linear normal form is written as
\begin{equation}
 R(s)=W(s)\,R_{\mathrm{rot}}(\mu_1,\mu_2)\,W^{-1}(s),
 \label{eq:normal-form}
\end{equation}
where
\begin{equation}
 R_{\mathrm{rot}}(\mu_1,\mu_2)
 =\diag\!\left(R_2(\mu_1),R_2(\mu_2)\right),
 \qquad
 R_2(\mu)=
 \begin{pmatrix}
  \cos\mu&\sin\mu\\
  -\sin\mu&\cos\mu
 \end{pmatrix},
\end{equation}
and
\begin{equation}
 W^T\mat S W=\mat S.
\end{equation}
The normal-mode phases at turn $N$ are
\begin{equation}
 \Phi_j(N)=2\pi Q_jN+\phi_{j0},
 \qquad j=1,2.
\end{equation}
Using actions $J_j$, the normalized mode coordinates are chosen as
\begin{equation}
 \widehat{\mat z}(N)=
 \begin{pmatrix}
  \sqrt{2J_1}\cos\Phi_1\\
  -\sqrt{2J_1}\sin\Phi_1\\
  \sqrt{2J_2}\cos\Phi_2\\
  -\sqrt{2J_2}\sin\Phi_2
 \end{pmatrix},
 \qquad
 \mat z(N)=W\widehat{\mat z}(N).
 \label{eq:mr-orbit-matrix}
\end{equation}

Define the complex eigenvectors
\begin{equation}
 \mat v_j=W_{:,2j-1}+\ii W_{:,2j}.
\end{equation}
Then the complete turn-by-turn orbit is
\begin{equation}
 \boxed{
 \mat z(N)=\sum_{j=1}^{2}\sqrt{2J_j}\,
 \Ree\!\left[\mat v_j\e^{\ii\Phi_j(N)}\right].
 }
 \label{eq:mr-orbit-complex}
\end{equation}
This form is independent of the phase gauge used for each eigenvector, provided the same gauge is used consistently in $W$ and in $\Phi_j$.

\subsection{Mais--Ripken lattice functions from $W$}

Partition $W$ into $2\times2$ blocks by physical plane and normal mode,
\begin{equation}
 W=
 \begin{pmatrix}
  W_{x1}&W_{x2}\\
  W_{y1}&W_{y2}
 \end{pmatrix}.
\end{equation}
For $q\in\{x,y\}$ and $j\in\{1,2\}$, define the projected mode matrix
\begin{equation}
 \Sigma_{qj}\equiv W_{qj}W_{qj}^T
 =
 \begin{pmatrix}
  \beta_{qj}&-\alpha_{qj}\\
  -\alpha_{qj}&\gamma_{qj}
 \end{pmatrix}.
 \label{eq:mr-sigma}
\end{equation}
Equivalently, if the two columns of $W_{qj}$ are denoted by $(w_{q,2j-1},w_{q,2j})$ and $(w_{p_q,2j-1},w_{p_q,2j})$, then
\begin{align}
 \beta_{qj}&=w_{q,2j-1}^2+w_{q,2j}^2,\\
 \alpha_{qj}&=-w_{q,2j-1}w_{p_q,2j-1}
               -w_{q,2j}w_{p_q,2j},\\
 \gamma_{qj}&=w_{p_q,2j-1}^2+w_{p_q,2j}^2.
 \label{eq:mr-functions-from-w}
\end{align}
Unlike an uncoupled Courant--Snyder triplet, an individual projected triplet need not satisfy $\beta\gamma-\alpha^2=1$.
These projected functions are the Mais--Ripken lattice functions~\cite{MaisRipken,LebedevBogacz}.

Writing
\begin{equation}
 v_{qj}=W_{q,2j-1}+\ii W_{q,2j}
       =\sqrt{\beta_{qj}}\e^{\ii\chi_{qj}},
\end{equation}
the position and canonical-momentum signals are
\begin{align}
 q(N)&=\sum_{j=1}^{2}\sqrt{2J_j\beta_{qj}}
       \cos\!\left(\Phi_j+\chi_{qj}\right),
 \label{eq:mr-position}\\
 p_q(N)&=\sum_{j=1}^{2}\sqrt{2J_j\gamma_{qj}}
       \cos\!\left(\Phi_j+\chi_{p_qj}\right),
 \label{eq:mr-momentum}
\end{align}
where $\chi_{p_qj}=\arg(W_{p_q,2j-1}+\ii W_{p_q,2j})$.  To avoid confusion
with the ET normalizing matrices $A_q$, let $\mathcal A_u(Q_j)$ denote the
amplitude of the Fourier line at tune $Q_j$ in a signal $u(N)$.  The position
and canonical-momentum line amplitudes are therefore
\begin{equation}
 \mathcal A_q(Q_j)=\sqrt{2J_j\beta_{qj}},
 \qquad
 \mathcal A_{p_q}(Q_j)=\sqrt{2J_j\gamma_{qj}}.
 \label{eq:mr-line-amplitudes}
\end{equation}

\subsection{Closest-tune-approach indicator}

A commonly used global coupling indicator can be obtained from the relative
amplitudes and phases of the two tune lines~\cite{luo2004phase,Jones:2020hhl}.
For the horizontal-like and vertical-like modes, respectively, define
\begin{equation}
 r_1\equiv
 \frac{\mathcal A_y(Q_1)}{\mathcal A_x(Q_1)}
 =\sqrt{\frac{\beta_{y1}}{\beta_{x1}}},
 \qquad
 r_2\equiv
 \frac{\mathcal A_x(Q_2)}{\mathcal A_y(Q_2)}
 =\sqrt{\frac{\beta_{x2}}{\beta_{y2}}},
\end{equation}
and the cross-plane phase differences
\begin{equation}
 \Delta\phi_1=\chi_{y1}-\chi_{x1},
 \qquad
 \Delta\phi_2=\chi_{x2}-\chi_{y2}.
\end{equation}
The closest-tune-approach magnitude is then~\cite{Jones:2020hhl}
\begin{equation}
 \left|C^-\right|
 =\frac{2\sqrt{r_1r_2}\,|Q_1-Q_2|}{1+r_1r_2}.
 \label{eq:cmin}
\end{equation}
Together with $\chi(s)=\Delta\phi_1(s)$, this defines the local complex
quantity
\begin{equation}
 C^-(s)=|C^-|\e^{\ii\chi(s)}.
\end{equation}
Only its phase depends on the observation point $s$.  In first-order
perturbation theory it is related to the skew-quadrupole distribution by
\begin{equation}
 C^-=\frac{1}{2\pi}\int_0^L
 \sqrt{\beta_x\beta_y}\,k_s
 \e^{\ii[\Phi_x-\Phi_y-2\pi\Delta s/L]}\,\dee l,
\end{equation}
where $\Delta$ is the difference of the uncoupled fractional tunes.  For a
more robust scalar estimate, \cref{eq:cmin} can be evaluated at every
observation point and averaged around the ring~\cite{tobias_coupling_note}.

\section{Edwards--Teng decoupling and modal Twiss parameters}
\label{sec:et-twiss}

\subsection{Decoupled modal maps}

The Edwards--Teng construction decouples the physical four-dimensional map
into two stable two-dimensional modal maps~\cite{EdwardsTeng,laurent_coupling_1,laurent_coupling_2}.
Following the notation used throughout this manual,
partition the one-turn map as
\begin{equation}
 R=\begin{pmatrix}A&B\\C&D\end{pmatrix},
\end{equation}
and denote the horizontal-like and vertical-like decoupled maps by $E$ and
$F$, respectively.  The purpose of the Edwards--Teng transformation is to
find a symplectic matrix $T_{\mathrm{ET}}$ such that
\begin{equation}
 T_{\mathrm{ET}}^{-1}RT_{\mathrm{ET}}
 =\begin{pmatrix}E&0\\0&F\end{pmatrix},
 \label{eq:et-decoupled-map}
\end{equation}
or, equivalently,
\begin{equation}
 R=T_{\mathrm{ET}}
 \begin{pmatrix}E&0\\0&F\end{pmatrix}
 T_{\mathrm{ET}}^{-1}.
 \label{eq:et-map-reconstruction}
\end{equation}
Thus $E$ and $F$ are the diagonal blocks of the transformed, decoupled map;
in a coupled lattice they are not the physical diagonal blocks $A$ and $D$ of
$R$.

Here and below,
\begin{equation}
 I_2=\begin{pmatrix}1&0\\0&1\end{pmatrix}
\end{equation}
denotes the $2\times2$ identity matrix.  In the Edwards--Teng convention used
here~\cite{EdwardsTeng}, the decoupling matrix is parametrized as
\begin{equation}
 T_{\mathrm{ET}}
 =g
 \begin{pmatrix}
  I_2&\bar R_{\mathrm{ET}}\\
  -R_{\mathrm{ET}}&I_2
 \end{pmatrix},
 \qquad
 g=\frac{1}{\sqrt{1+\det R_{\mathrm{ET}}}}>0.
 \label{eq:tet}
\end{equation}
We now derive $R_{\mathrm{ET}}$ from the physical map.  Extend the symplectic
conjugate to the full four-dimensional map.  Since $R$ is symplectic,
\begin{equation}
 \overline R\equiv-\mat S R^T\mat S=R^{-1}.
\end{equation}
For any real $2\times2$ matrix $Z$,
\begin{equation}
 Z+\bar Z=(\Tr Z)I_2,
 \qquad
 \bar Z Z=Z\bar Z=(\det Z)I_2.
 \label{eq:two-by-two-conjugate-identities}
\end{equation}
The reason for considering $R+\overline R$ is that it exposes the ET modal
planes.  Indeed, using \cref{eq:et-map-reconstruction} and
$\overline R=R^{-1}$ gives
\begin{equation}
 R+\overline R
 =T_{\mathrm{ET}}
  \begin{pmatrix}
   E+E^{-1}&0\\0&F+F^{-1}
  \end{pmatrix}
  T_{\mathrm{ET}}^{-1}.
 \label{eq:R-plus-Rbar-et-basis}
\end{equation}
Because $E$ and $F$ are $2\times2$ symplectic matrices,
\cref{eq:two-by-two-conjugate-identities} implies
\begin{equation}
 E+E^{-1}=(\Tr E)I_2,
 \qquad
 F+F^{-1}=(\Tr F)I_2.
 \label{eq:modal-map-plus-inverse}
\end{equation}
Consequently, provided $\Tr E\ne\Tr F$, the two eigenspaces of
$R+\overline R$ are precisely the two ET modal planes.  The first block
column of \cref{eq:tet} shows that every vector in the horizontal-like plane
has the form
\begin{equation}
 T_{\mathrm{ET}}
 \begin{pmatrix}\mat x\\0\end{pmatrix}
 =g\begin{pmatrix}\mat x\\-R_{\mathrm{ET}}\mat x\end{pmatrix}.
 \label{eq:first-et-plane-from-tet}
\end{equation}
Thus, if an eigenvector of $R+\overline R$ in this plane is written as
$(\mat x,\mat y)^T$, its two blocks satisfy
\begin{equation}
 \mat y=-R_{\mathrm{ET}}\mat x.
 \label{eq:ret-as-modal-plane-slope}
\end{equation}
We can therefore obtain $R_{\mathrm{ET}}$ by calculating the same eigenspace
directly from the physical blocks of $R$ and comparing its $\mat y$ and
$\mat x$ components.

Therefore, on defining
\begin{equation}
 a\equiv\Tr A,
 \qquad d\equiv\Tr D,
 \qquad X\equiv C+\bar B,
 \label{eq:et-a-d-X}
\end{equation}
one obtains
\begin{equation}
 R+\overline R
 =\begin{pmatrix}
   A+\bar A&B+\bar C\\
   C+\bar B&D+\bar D
  \end{pmatrix}
 =\begin{pmatrix}
   aI_2&\bar X\\X&dI_2
  \end{pmatrix}.
 \label{eq:R-plus-Rbar}
\end{equation}

Let $\Lambda$ be an eigenvalue of $R+\overline R$ and write a corresponding
two-block eigenvector as $(\mat x,\mat y)^T$.  Equation~\eqref{eq:R-plus-Rbar}
gives
\begin{align}
 (a-\Lambda)\mat x+\bar X\mat y&=0,
 \label{eq:et-eigenvector-upper}\\
 X\mat x+(d-\Lambda)\mat y&=0.
 \label{eq:et-eigenvector-lower}
\end{align}
For $\Lambda\ne d$, the second equation gives
\begin{equation}
 \mat y=\frac{X}{\Lambda-d}\mat x.
 \label{eq:y-from-x-et}
\end{equation}
Substitution into the first equation, together with
$\bar X X=(\det X)I_2$, yields the scalar characteristic equation
\begin{equation}
 (a-\Lambda)(d-\Lambda)-\det X=0.
 \label{eq:et-lambda-characteristic}
\end{equation}
Expanding the product makes this an ordinary quadratic equation in
$\Lambda$:
\begin{equation}
 \Lambda^2-(a+d)\Lambda+ad-\det X=0.
 \label{eq:et-lambda-quadratic}
\end{equation}
The quadratic formula therefore gives
\begin{align}
 \Lambda_{\pm}
 &=\frac{a+d\mathbin{\pm}
  \sqrt{(a+d)^2-4(ad-\det X)}}{2} \notag\\
 &=\frac{a+d}{2}\mathbin{\pm}\sqrt{\Delta},
 \qquad
 \Delta\equiv\frac14(a-d)^2+\det X,
 \label{eq:et-delta}
\end{align}
where the second line follows because
\begin{equation}
 (a+d)^2-4(ad-\det X)
 =(a-d)^2+4\det X
 =4\Delta.
 \label{eq:et-discriminant}
\end{equation}

The horizontal-like root must approach $a=\Tr A$ when the coupling block
$X$ tends to zero.  In that limit $\det X\to0$ and hence
\begin{equation}
 \sqrt{\Delta}\longrightarrow\frac{|a-d|}{2}.
 \label{eq:et-uncoupled-discriminant}
\end{equation}
Consequently, the sign that selects the horizontal-like root depends on
the ordering of $a$ and $d$:
\begin{equation}
 \begin{aligned}
  a>d:&\qquad \Lambda_+\longrightarrow a,
                    &\Lambda_-\longrightarrow d,\\
  a<d:&\qquad \Lambda_-\longrightarrow a,
                    &\Lambda_+\longrightarrow d.
 \end{aligned}
 \label{eq:et-uncoupled-root-selection}
\end{equation}
Both cases are encoded by $\operatorname{sign}(a-d)$.  Thus, for $a\ne d$,
\begin{equation}
 \begin{aligned}
  \Lambda_1
  &=\frac{a+d}{2}
    +\operatorname{sign}(a-d)\sqrt{\Delta},
  &\Lambda_1-d
  &=\frac{a-d}{2}
    +\operatorname{sign}(a-d)\sqrt{\Delta},\\
  \Lambda_2
  &=\frac{a+d}{2}
    -\operatorname{sign}(a-d)\sqrt{\Delta},
  &\Lambda_2-a
  &=\frac{d-a}{2}
    -\operatorname{sign}(a-d)\sqrt{\Delta}.
 \end{aligned}
 \label{eq:horizontal-et-root}
\end{equation}
The first modal subspace in $T_{\mathrm{ET}}$ consists of vectors
$(\mat x,-R_{\mathrm{ET}}\mat x)^T$.  Comparing this with
\cref{eq:y-from-x-et} for $\Lambda=\Lambda_1$ therefore gives
\begin{equation}
 \boxed{
 R_{\mathrm{ET}}
 =-\left[
  \frac12\left(\Tr A-\Tr D\right)
  +\operatorname{sign}\!\left(\Tr A-\Tr D\right)\sqrt{\Delta}
 \right]^{-1}
 \left(C+\bar B\right).
 }
 \label{eq:ret-from-map-blocks}
\end{equation}

This eigenvector construction also proves the block diagonalization.  Indeed,
$R$ commutes with $R+R^{-1}$, so each nondegenerate eigenspace of
$R+R^{-1}$ is invariant under $R$.  Writing $K=R_{\mathrm{ET}}$, invariance
of the first subspace gives
\begin{equation}
 R\begin{pmatrix}I_2\\-K\end{pmatrix}
 =\begin{pmatrix}I_2\\-K\end{pmatrix}E,
 \qquad
 E=A-BK,
 \label{eq:first-et-invariant-subspace}
\end{equation}
whose lower block is the matrix Riccati equation
\begin{equation}
 C+KA-DK-KBK=0.
 \label{eq:et-riccati}
\end{equation}
The second modal eigenspace similarly gives
\begin{equation}
 R\begin{pmatrix}\bar K\\I_2\end{pmatrix}
 =\begin{pmatrix}\bar K\\I_2\end{pmatrix}F,
 \qquad
 F=D+C\bar K=D+KB.
 \label{eq:second-et-invariant-subspace}
\end{equation}
The equality $C\bar K=KB$ follows from
$K=-X/(\Lambda_1-d)$ and $\det B=\det C$.  The latter is itself an immediate
consequence of the symplectic block identities
$A\bar A+B\bar B=I_2=A\bar A+C\bar C$ and
\cref{eq:two-by-two-conjugate-identities}.
The two block columns in these equations are precisely the two block columns
of $T_{\mathrm{ET}}/g$, so together they imply
\cref{eq:et-decoupled-map}.  Moreover,
\begin{equation}
 \left(\frac{T_{\mathrm{ET}}}{g}\right)^T
 \mat S
 \left(\frac{T_{\mathrm{ET}}}{g}\right)
 =\left(1+\det K\right)\mat S,
\end{equation}
which explains the normalization $g=(1+\det K)^{-1/2}$ in
\cref{eq:tet}.  At a degeneracy the separation into the two modal subspaces
is not unique, and the branch must instead be selected by continuity from a
neighboring lattice point.
By construction, $T_{\mathrm{ET}}$ maps modal coordinates
$\mat z_{\mathrm{ET}}$ to physical coordinates $\mat z$:
\begin{equation}
 \mat z=T_{\mathrm{ET}}\mat z_{\mathrm{ET}}.
\end{equation}

Write the two stable modal maps as
\begin{equation}
 E=A_xR_2(\mu_1)A_x^{-1},
 \qquad
 F=A_yR_2(\mu_2)A_y^{-1},
 \label{eq:et-modal-map}
\end{equation}
where $R_2$ is the rotation matrix defined above and
\begin{equation}
 A_q=
 \begin{pmatrix}
  \sqrt{\beta_q}&0\\
  -\alpha_q/\sqrt{\beta_q}&1/\sqrt{\beta_q}
 \end{pmatrix}.
 \label{eq:Aq}
\end{equation}
For $q\in\{x,y\}$, set $\gamma_q=(1+\alpha_q^2)/\beta_q$.  The entries of
the horizontal-like map are then
\begin{equation}
 E=
 \begin{pmatrix}
  E_{11}&E_{12}\\E_{21}&E_{22}
 \end{pmatrix}
 =
 \begin{pmatrix}
  \cos\mu_1+\alpha_x\sin\mu_1&\beta_x\sin\mu_1\\
  -\gamma_x\sin\mu_1&\cos\mu_1-\alpha_x\sin\mu_1
 \end{pmatrix}.
 \label{eq:E-twiss-form}
\end{equation}
Consequently, its phase advance and Twiss parameters can be recovered
directly from $E$:
\begin{align}
 \cos\mu_1
 &=\frac12\Tr E,
 \label{eq:cosmu1-from-E}\\
 \sin\mu_1
 &=\operatorname{sign}(E_{12})
   \sqrt{-E_{12}E_{21}
   -\left(\frac{E_{11}-E_{22}}{2}\right)^2},
 \label{eq:sinmu1-from-E}\\
 \beta_x&=\frac{E_{12}}{\sin\mu_1},
 &\gamma_x&=-\frac{E_{21}}{\sin\mu_1},
 &\alpha_x&=\frac{E_{11}-E_{22}}{2\sin\mu_1}.
 \label{eq:et-twiss-from-E}
\end{align}
The vertical-like quantities follow from the same equations with
$E\to F$, $(\beta_x,\alpha_x,\gamma_x)\to
(\beta_y,\alpha_y,\gamma_y)$, and $\mu_1\to\mu_2$.

Following the propagation subsection of the Xsuite physics manual, the
Edwards--Teng Twiss parameters are denoted by $(\beta_x,\alpha_x)$ for the
horizontal-like mode and $(\beta_y,\alpha_y)$ for the vertical-like mode.
They are not the projected Mais--Ripken functions
$(\beta_{x1},\alpha_{x1})$ and $(\beta_{y2},\alpha_{y2})$.

Its inverse is
\begin{equation}
 A_q^{-1}=
 \begin{pmatrix}
  \beta_q^{-1/2}&0\\
  \alpha_q\beta_q^{-1/2}&\beta_q^{1/2}
 \end{pmatrix}
 \label{eq:Aq-inverse}
\end{equation}
and maps physical modal coordinates in plane $q$ to Edwards--Teng normalized
coordinates:
\begin{equation}
 \begin{pmatrix}\widetilde q\\\widetilde p_q\end{pmatrix}
 =A_q^{-1}\begin{pmatrix}q\\p_q\end{pmatrix}.
 \label{eq:cs-normalization}
\end{equation}
In the uncoupled limit, these modal Twiss parameters reduce to the usual
horizontal and vertical Courant--Snyder Twiss parameters.

\subsection{Relation between $W$ and the Edwards--Teng transformation}

The matrix $A_q^{-1}$ maps modal coordinates to normalized coordinates,
whereas $A_q$ performs the inverse operation.  Throughout this subsection and
all that follows, we assume that the phases of the two transverse eigenvectors
of $W$ have already been chosen such that $v_{1,x}$ and $v_{2,y}$ are real and
positive.  With
\begin{equation}
 \mat v_1=W_{:,1}+\ii W_{:,2},
 \qquad
 \mat v_2=W_{:,3}+\ii W_{:,4},
\end{equation}
this means explicitly
\begin{equation}
 \Imm(v_{1,x})=0,
 \qquad
 \Ree(v_{1,x})>0,
 \qquad
 \Imm(v_{2,y})=0,
 \qquad
 \Ree(v_{2,y})>0.
 \label{eq:et-rephasing-of-w}
\end{equation}
This is the Edwards--Teng phase gauge used in the implementation.

We shall now prove that this phase choice makes the normalized coordinates
associated with $W$ and with the ET construction identical.  Initially denote
them separately by
\begin{equation}
 \widehat{\mat z}_W=W^{-1}\mat z,
 \qquad
 \widehat{\mat z}_{\mathrm{ET}}
 =\diag(A_x^{-1},A_y^{-1})T_{\mathrm{ET}}^{-1}\mat z.
 \label{eq:two-normalized-coordinates}
\end{equation}
Define
\begin{equation}
 \mathcal A\equiv\diag(A_x,A_y),
 \qquad
 U\equiv T_{\mathrm{ET}}\mathcal A.
 \label{eq:U-definition}
\end{equation}
The second definition in \cref{eq:two-normalized-coordinates} is then
equivalent to
\begin{equation}
 \mat z=U\widehat{\mat z}_{\mathrm{ET}}.
 \label{eq:z-from-et-normalized}
\end{equation}

Both $U$ and $W$ reduce the physical map to the same block rotation.  To see
this explicitly, use
$U^{-1}=\mathcal A^{-1}T_{\mathrm{ET}}^{-1}$ and substitute the ET
block-diagonalization from \cref{eq:et-decoupled-map}:
\begin{align}
 U^{-1}RU
 &=\mathcal A^{-1}T_{\mathrm{ET}}^{-1}
   RT_{\mathrm{ET}}\mathcal A
 \nonumber\\
 &=\mathcal A^{-1}\diag(E,F)\mathcal A
 \nonumber\\
 &=\diag(A_x^{-1}EA_x,A_y^{-1}FA_y)
 \nonumber\\
 &=\diag(R_2(\mu_1),R_2(\mu_2))
 \nonumber\\
 &=R_{\mathrm{rot}}(\mu_1,\mu_2).
 \label{eq:U-normalizes-R}
\end{align}
The penultimate equality follows directly from
$E=A_xR_2(\mu_1)A_x^{-1}$ and $F=A_yR_2(\mu_2)A_y^{-1}$ in
\cref{eq:et-modal-map}.
Likewise, \cref{eq:normal-form} gives
\begin{equation}
 W^{-1}RW=R_{\mathrm{rot}}(\mu_1,\mu_2).
\end{equation}
Define the change of normalized basis
\begin{equation}
 D\equiv U^{-1}W
 \label{eq:D-definition}
\end{equation}
so that $W=UD$.  Substituting this relation into the preceding normal-form
equation gives
\begin{align}
 R_{\mathrm{rot}}(\mu_1,\mu_2)
 &=W^{-1}RW \nonumber\\
 &=(UD)^{-1}R(UD) \nonumber\\
 &=D^{-1}(U^{-1}RU)D \nonumber\\
 &=D^{-1}R_{\mathrm{rot}}(\mu_1,\mu_2)D.
\end{align}
Multiplication from the left by $D$ therefore shows explicitly that $D$
commutes with the block rotation:
\begin{equation}
 DR_{\mathrm{rot}}(\mu_1,\mu_2)
 =R_{\mathrm{rot}}(\mu_1,\mu_2)D.
\end{equation}
For distinct, nondegenerate tunes and a fixed ordering of the two modes, a
symplectic matrix with this property consists of an independent phase rotation
in each mode:
\begin{equation}
 D=\diag\!\left(R_2(\delta_1),R_2(\delta_2)\right).
 \label{eq:D-phase-rotations}
\end{equation}
It follows from $W=UD$ that the two normalized vectors are related by
\begin{equation}
 \widehat{\mat z}_{\mathrm{ET}}=D\widehat{\mat z}_W.
 \label{eq:normalized-coordinate-gauges}
\end{equation}
It remains to use the phase convention in \cref{eq:et-rephasing-of-w} to
determine $D$.

Substituting the explicit form of $T_{\mathrm{ET}}$ from \cref{eq:tet} into
the definition of $U$ in \cref{eq:U-definition} shows that its upper-left and
lower-right blocks are $gA_x$ and $gA_y$, respectively.
Because the first row of each $A_q$ is $(\sqrt{\beta_q},0)$, the corresponding
complex position components of $U$ are the positive real numbers
$g\sqrt{\beta_x}$ and $g\sqrt{\beta_y}$.  Right multiplication of a real
column pair by $R_2(\delta)$ multiplies its associated complex vector by
$e^{\ii\delta}$.  Since $W=UD$, it follows that
\begin{equation}
 v_{1,x}=g\sqrt{\beta_x}\,e^{\ii\delta_1},
 \qquad
 v_{2,y}=g\sqrt{\beta_y}\,e^{\ii\delta_2}.
 \label{eq:principal-components-with-D}
\end{equation}
The factors $g\sqrt{\beta_x}$ and $g\sqrt{\beta_y}$ are positive.  Since the
left-hand sides are also real and positive by \cref{eq:et-rephasing-of-w},
\begin{equation}
 e^{\ii\delta_1}=e^{\ii\delta_2}=1,
 \qquad
 D=\diag(I_2,I_2).
 \label{eq:et-gauge-fixes-D}
\end{equation}
Therefore $W=UD=U$, and \cref{eq:normalized-coordinate-gauges} gives
\begin{equation}
 \boxed{
 \widehat{\mat z}_{\mathrm{ET}}=\widehat{\mat z}_W
 \equiv\widehat{\mat z}.
 }
 \label{eq:equal-normalized-coordinates}
\end{equation}
Finally, $W=U$ and \cref{eq:U-definition} give
\begin{equation}
 \boxed{
 W=T_{\mathrm{ET}}\diag(A_x,A_y).
 }
 \label{eq:w-from-tet}
\end{equation}
All equations below that involve the individual columns or blocks of $W$ use
this explicitly rephased matrix.
Substituting the explicit ET matrix from \cref{eq:tet} into
\cref{eq:w-from-tet} first gives
\begin{equation}
 \begin{aligned}
 W
 &=g
 \begin{pmatrix}
  I_2&\bar R_{\mathrm{ET}}\\
  -R_{\mathrm{ET}}&I_2
 \end{pmatrix}
 \begin{pmatrix}A_x&0\\0&A_y\end{pmatrix}\\
 &=\begin{pmatrix}
  g A_x&g\bar R_{\mathrm{ET}}A_y\\
  -g R_{\mathrm{ET}}A_x&g A_y
 \end{pmatrix}.
 \end{aligned}
 \label{eq:w-et-block-product}
\end{equation}
To factor the modal normalizers $A_x$ and $A_y$ from the left, define the
ET-normalized coupling matrix
\begin{equation}
 \boxed{
 \mathcal C=g A_x^{-1}\bar R_{\mathrm{ET}}A_y.
 }
 \label{eq:c-from-ret}
\end{equation}
The upper-right block in \cref{eq:w-et-block-product} is consequently
\begin{equation}
 A_x\mathcal C=g\bar R_{\mathrm{ET}}A_y.
\end{equation}
For the lower-left block, use
\begin{equation}
 \overline{PQ}=\bar Q\,\bar P,
 \qquad
 \overline{\bar R_{\mathrm{ET}}}=R_{\mathrm{ET}},
 \qquad
 \bar A_q=A_q^{-1},
\end{equation}
where the last identity follows because $A_q$ is symplectic.  Taking the
symplectic conjugate of \cref{eq:c-from-ret} gives
\begin{align}
 \bar{\mathcal C}
 &=g
 \overline{A_x^{-1}\bar R_{\mathrm{ET}}A_y}
 \nonumber\\
 &=g A_y^{-1}R_{\mathrm{ET}}A_x,
\end{align}
and therefore
\begin{equation}
 A_y\bar{\mathcal C}=g R_{\mathrm{ET}}A_x.
\end{equation}
Using these two block identities, \cref{eq:w-et-block-product} becomes
\begin{equation}
 \boxed{
 W=
 \begin{pmatrix}A_x&0\\0&A_y\end{pmatrix}
 \begin{pmatrix}
  g I_2&\mathcal C\\
  -\bar{\mathcal C}&g I_2
 \end{pmatrix}.
 }
 \label{eq:w-factorization}
\end{equation}
Symplecticity requires
\begin{equation}
 g^2+\det\mathcal C=1.
 \label{eq:g-c-condition}
\end{equation}
Equivalently, if $W$ is partitioned as in \cref{eq:w-factorization}, its
phase-gauged blocks obey
\begin{equation}
 \boxed{
 A_x^{-1}W_{x1}=g I_2,\qquad
 A_x^{-1}W_{x2}=\mathcal C,\qquad
 A_y^{-1}W_{y1}=-\bar{\mathcal C},\qquad
 A_y^{-1}W_{y2}=g I_2.
 }
 \label{eq:w-blocks-normalized}
\end{equation}
For later use, define
\begin{equation}
 u\equiv\det\mathcal C
 =\frac{\det R_{\mathrm{ET}}}{1+\det R_{\mathrm{ET}}}
 =1-g^2.
 \label{eq:u-def}
\end{equation}
For the ordinary trigonometric ET branch, $u=\sin^2\phi$ and
$g^2=\cos^2\phi$, where $\phi$ is the Teng rotation angle.

\subsection{Relation between Edwards--Teng and Mais--Ripken functions}

Using the definition $\Sigma_{qj}=W_{qj}W_{qj}^T$ from
\cref{eq:mr-sigma} and substituting the four blocks of $W$ from
\cref{eq:w-factorization} gives
\begin{align}
 \Sigma_{x1}&=g^2 A_xA_x^T,
 &\Sigma_{x2}&=A_x\mathcal C\mathcal C^TA_x^T,
 \label{eq:sigmax-factorization}\\
 \Sigma_{y1}&=A_y\bar{\mathcal C}(\bar{\mathcal C})^TA_y^T,
 &\Sigma_{y2}&=g^2 A_yA_y^T.
 \label{eq:sigmay-factorization}
\end{align}
The principal Mais--Ripken functions are therefore
\begin{align}
 \beta_{x1}&=g^2\beta_x,
 &\alpha_{x1}&=g^2\alpha_x,
 &\gamma_{x1}&=g^2\gamma_x,
 \label{eq:mr-et-main-x}\\
 \beta_{y2}&=g^2\beta_y,
 &\alpha_{y2}&=g^2\alpha_y,
 &\gamma_{y2}&=g^2\gamma_y.
 \label{eq:mr-et-main-y}
\end{align}
The determinant of either principal projected block is therefore
\begin{equation}
 \beta_{x1}\gamma_{x1}-\alpha_{x1}^2
 =\beta_{y2}\gamma_{y2}-\alpha_{y2}^2
 =g^4.
 \label{eq:main-determinants}
\end{equation}
Consequently, $g$ is obtained directly from the principal Mais--Ripken functions:
\begin{equation}
 \boxed{
 g^2
 =\sqrt{\beta_{x1}\gamma_{x1}-\alpha_{x1}^2}
 =\sqrt{\beta_{y2}\gamma_{y2}-\alpha_{y2}^2}.
 }
 \label{eq:g-from-mr}
\end{equation}
The ET modal Twiss parameters then follow by division by $g^2$:
\begin{equation}
 \boxed{
 \beta_x=\frac{\beta_{x1}}{g^2},\quad
 \alpha_x=\frac{\alpha_{x1}}{g^2},\qquad
 \beta_y=\frac{\beta_{y2}}{g^2},\quad
 \alpha_y=\frac{\alpha_{y2}}{g^2}.
 }
 \label{eq:et-from-main-mr}
\end{equation}

To make the remaining Mais--Ripken functions equally explicit, write
\begin{equation}
 \mathcal C=
 \begin{pmatrix}
  c_{11}&c_{12}\\c_{21}&c_{22}
 \end{pmatrix}.
 \label{eq:c-components}
\end{equation}
Then
\begin{align}
 \beta_{x2}
 &=\beta_x(c_{11}^2+c_{12}^2),
 \label{eq:betax2-et}\\
 \alpha_{x2}
 &=\alpha_x(c_{11}^2+c_{12}^2)
   -(c_{11}c_{21}+c_{12}c_{22}),
 \label{eq:alphax2-et}\\
 \gamma_{x2}
 &=\frac{
  (c_{21}-\alpha_xc_{11})^2
  +(c_{22}-\alpha_xc_{12})^2
 }{\beta_x},
 \label{eq:gammax2-et}
\end{align}
and, since
\begin{equation}
 -\bar{\mathcal C}
 =
 \begin{pmatrix}
  -c_{22}&c_{12}\\
  c_{21}&-c_{11}
 \end{pmatrix},
\end{equation}
one also has
\begin{align}
 \beta_{y1}
 &=\beta_y(c_{22}^2+c_{12}^2),
 \label{eq:betay1-et}\\
 \alpha_{y1}
 &=\alpha_y(c_{22}^2+c_{12}^2)
   +c_{21}c_{22}+c_{11}c_{12},
 \label{eq:alphay1-et}\\
 \gamma_{y1}
 &=\frac{
  (c_{21}+\alpha_yc_{22})^2
  +(c_{11}+\alpha_yc_{12})^2
 }{\beta_y}.
 \label{eq:gammay1-et}
\end{align}
The cross-plane projected determinants satisfy
\begin{equation}
 \begin{aligned}
 \beta_{x2}\gamma_{x2}-\alpha_{x2}^2
 &=\beta_{y1}\gamma_{y1}-\alpha_{y1}^2\\
 &=(1-g^2)^2.
 \end{aligned}
 \label{eq:cross-determinants}
\end{equation}

\subsection{Propagation of Edwards--Teng parameters}
\label{sec:et-propagation}

Let the transfer map from location $a$ to location $b$ be partitioned as
\begin{equation}
 R^{ab}=
 \begin{pmatrix}
  A^{ab}&B^{ab}\\C^{ab}&D^{ab}
 \end{pmatrix}.
 \label{eq:Rab-blocks}
\end{equation}
Given the ET quantities $E^a$, $F^a$, and $R_{\mathrm{ET}}^a$ at $a$, define
the two effective modal transfer matrices
\begin{align}
 E^{ab}&=A^{ab}-B^{ab}R_{\mathrm{ET}}^a,
 \label{eq:Eab-definition}\\
 F^{ab}&=D^{ab}+C^{ab}\bar R_{\mathrm{ET}}^a.
 \label{eq:Fab-definition}
\end{align}
The decoupled one-turn maps and the ET coupling matrix at $b$ are then
\begin{align}
 E^b
 &=E^{ab}E^a\bar E^{ab}/\det E^{ab},
 \label{eq:E-propagation}\\
 F^b
 &=F^{ab}F^a\bar F^{ab}/\det F^{ab},
 \label{eq:F-propagation}\\
 R_{\mathrm{ET}}^b
 &=-\left(C^{ab}-D^{ab}R_{\mathrm{ET}}^a\right)
   \bar E^{ab}/\det E^{ab}.
 \label{eq:ret-propagation}
\end{align}
The factors $\bar E^{ab}/\det E^{ab}$ and
$\bar F^{ab}/\det F^{ab}$ are simply $(E^{ab})^{-1}$ and
$(F^{ab})^{-1}$, respectively.

Writing $E^{ab}_{ij}$ for the entries of $E^{ab}$, the horizontal ET Twiss
parameters and phase propagate according to
\begin{align}
 \beta_x^b
 &=\frac{
   \left(E^{ab}_{11}\beta_x^a-E^{ab}_{12}\alpha_x^a\right)^2
   +\left(E^{ab}_{12}\right)^2
  }{\det(E^{ab})\,\beta_x^a},
 \label{eq:et-beta-propagation}\\
 \alpha_x^b
 &=-\frac{1}{\det(E^{ab})\,\beta_x^a}
 \Big[
  \left(E^{ab}_{21}\beta_x^a-E^{ab}_{22}\alpha_x^a\right)
  \left(E^{ab}_{11}\beta_x^a-E^{ab}_{12}\alpha_x^a\right)
  +E^{ab}_{12}E^{ab}_{22}
 \Big],
 \label{eq:et-alpha-propagation}\\
 \mu_1^b
 &=\mu_1^a+operatorname{atan2}\!\left(
  E^{ab}_{12},
  E^{ab}_{11}\beta_x^a-E^{ab}_{12}\alpha_x^a
 \right),
 \label{eq:et-mu-propagation}\\
 \gamma_x^b&=\frac{1+(\alpha_x^b)^2}{\beta_x^b}.
 \label{eq:et-gamma-propagation}
\end{align}
The vertical formulas follow by replacing $E^{ab}$ with $F^{ab}$,
$x$ with $y$, and mode 1 with mode 2.

\section{Turn-by-turn motion in the RDT formalism}
\label{sec:rdt}

\subsection{Complex Courant--Snyder coordinates}

Using the Edwards--Teng normalized coordinates defined in
\cref{eq:cs-normalization}, the complex Courant--Snyder variables are
\begin{equation}
 h_{q,\pm}=\widetilde q\pm\ii\widetilde p_q
 =\sqrt{2I_q}\e^{\mp\ii\Phi_q}.
 \label{eq:complex-cs}
\end{equation}
Thus $h_{q,-}=\widetilde q-\ii\widetilde p_q$ and
\begin{equation}
 \widetilde q=\Ree(h_{q,-}),
 \qquad
 \widetilde p_q=-\Imm(h_{q,-}).
\end{equation}

The normal-form generating function is expanded as
\begin{equation}
 F=\sum_{jklm}f_{jklm}
 h_{x,+}^{j}h_{x,-}^{k}h_{y,+}^{l}h_{y,-}^{m}.
\end{equation}
To first order in the RDTs, the turn-by-turn motion is
\begin{align}
 h_{x,-}(s,N)
 &=\sqrt{2I_x}\e^{\ii\Phi_x}
 -2\ii\sum_{jklm}j f_{jklm}^{(s)}
 (2I_x)^{(j+k-1)/2}(2I_y)^{(l+m)/2}
 \nonumber\\[-0.2em]
 &\hspace{3.2cm}\times
 \e^{\ii[(1-j+k)\Phi_x+(m-l)\Phi_y]},
 \label{eq:general-hx-rdt}\\
 h_{y,-}(s,N)
 &=\sqrt{2I_y}\e^{\ii\Phi_y}
 -2\ii\sum_{jklm}l f_{jklm}^{(s)}
 (2I_x)^{(j+k)/2}(2I_y)^{(l+m-1)/2}
 \nonumber\\[-0.2em]
 &\hspace{3.2cm}\times
 \e^{\ii[(k-j)\Phi_x+(1-l+m)\Phi_y]}.
 \label{eq:general-hy-rdt}
\end{align}
Here $\Phi_x=2\pi Q_xN+\psi_{x,s,0}$ and similarly for $y$.

At first order in lattice perturbation theory, the local RDT at observation point $s$ can be written as
\begin{equation}
 f_{jklm}(s)=
 \frac{
  \displaystyle\sum_{w}h_{w,jklm}
  \e^{\ii[(j-k)\Delta\phi_{x,w\to s}+(l-m)\Delta\phi_{y,w\to s}]}
 }{
  1-\e^{2\pi\ii[(j-k)Q_x+(l-m)Q_y]}
 }.
 \label{eq:general-local-rdt}
\end{equation}
The coefficient $h_{w,jklm}$ is the corresponding Hamiltonian coefficient of source $w$.  For a thin skew quadrupole, $h_{w,1001}$ and $h_{w,1010}$ have the same local strength factor, proportional to
\[
 K_{1s,w}\sqrt{\beta_{x,w}\beta_{y,w}}/4;
\]
their different global behavior comes from their phase factors and resonance denominators in \cref{eq:rdt_def}.

In particular,
\begin{align}
 f_{1001}(s)&=
 \frac{\displaystyle\sum_w h_{w,1001}
 \e^{\ii(\Delta\phi_{x,w\to s}-\Delta\phi_{y,w\to s})}}
 {1-\e^{2\pi\ii(Q_x-Q_y)}},
 \label{eq:f1001-source}\\
 f_{1010}(s)&=
 \frac{\displaystyle\sum_w h_{w,1010}
 \e^{\ii(\Delta\phi_{x,w\to s}+\Delta\phi_{y,w\to s})}}
 {1-\e^{2\pi\ii(Q_x+Q_y)}}.
 \label{eq:f1010-source}
\end{align}

\subsection{The four linear-coupling monomials}

Reality of the Hamiltonian implies
\begin{equation}
 f_{jklm}^{*}=f_{kjml}.
\end{equation}
Hence
\begin{equation}
 \boxed{
 f_{0110}=f_{1001}^{*},
 \qquad
 f_{0101}=f_{1010}^{*}.
 }
 \label{eq:rdt-conjugates}
\end{equation}
Only terms with $j>0$ enter $h_{x,-}$, whereas only terms with $l>0$ enter
$h_{y,-}$.  Equations~\eqref{eq:general-hx-rdt} and
\eqref{eq:general-hy-rdt} give
\begin{align}
 \boxed{
 h_{x,-}=\sqrt{2I_x}\e^{\ii\Phi_x}
 -2\ii\sqrt{2I_y}
 \left(\fminus\e^{\ii\Phi_y}+\fplus\e^{-\ii\Phi_y}\right)
 }
 \label{eq:hx-linear-coupling}\\
 \boxed{
 h_{y,-}=\sqrt{2I_y}\e^{\ii\Phi_y}
 -2\ii\sqrt{2I_x}
 \left(\fminus^{*}\e^{\ii\Phi_x}+\fplus\e^{-\ii\Phi_x}\right).
 }
 \label{eq:hy-linear-coupling}
\end{align}
The difference-resonance term $f_{1001}$ produces the opposite-plane line with the same sign of frequency; the sum-resonance term $f_{1010}$ produces the negative-frequency contribution.  In a real position signal these two contributions interfere at the same observed tune line.

\section{Observable position and momentum amplitudes}
\label{sec:observables}

\subsection{Normalized-coordinate signals}

Using $\widetilde q=\Ree h_{q,-}$ and $\widetilde p_q=-\Imm h_{q,-}$, the horizontal signal becomes
\begin{align}
 \widetilde x(N)
 &=\sqrt{2I_x}\cos\Phi_x
 +2\sqrt{2I_y}\,
 \Imm\!\left[(\fminus-\fplus^{*})\e^{\ii\Phi_y}\right],
 \label{eq:xtilde-rdt}\\
 \widetilde p_x(N)
 &=-\sqrt{2I_x}\sin\Phi_x
 +2\sqrt{2I_y}\,
 \Ree\!\left[(\fminus+\fplus^{*})\e^{\ii\Phi_y}\right].
 \label{eq:pxtilde-rdt}
\end{align}
Similarly,
\begin{align}
 \widetilde y(N)
 &=\sqrt{2I_y}\cos\Phi_y
 +2\sqrt{2I_x}\,
 \Imm\!\left[(\fminus^{*}-\fplus^{*})\e^{\ii\Phi_x}\right],
 \label{eq:ytilde-rdt}\\
 \widetilde p_y(N)
 &=-\sqrt{2I_y}\sin\Phi_y
 +2\sqrt{2I_x}\,
 \Ree\!\left[(\fminus^{*}+\fplus^{*})\e^{\ii\Phi_x}\right].
 \label{eq:pytilde-rdt}
\end{align}
The four basic combinations are summarized in \cref{tab:rdt-combinations}.

\begin{table}[htbp]
\centering
\caption{RDT combinations controlling the cross-plane mode projection.}
\label{tab:rdt-combinations}
\begin{tabular}{@{}ll@{}}
\toprule
Observable & Cross-mode complex coefficient\\
\midrule
$x$ from mode 2 & $\fminus-\fplus^{*}$\\
$\widetilde p_x$ from mode 2 & $\fminus+\fplus^{*}$\\
$y$ from mode 1 & $(\fminus-\fplus)^{*}$\\
$\widetilde p_y$ from mode 1 & $(\fminus+\fplus)^{*}$\\
\bottomrule
\end{tabular}
\end{table}

For a particle carrying both mode actions, the measured position-line ratios are therefore
\begin{align}
 \frac{\mathcal A_x(Q_2)}{\mathcal A_x(Q_1)}
 &=2\sqrt{\frac{I_y}{I_x}}\,|\fminus-\fplus^{*}|,
 \label{eq:measured-x-rdt-ratio}\\
 \frac{\mathcal A_y(Q_1)}{\mathcal A_y(Q_2)}
 &=2\sqrt{\frac{I_x}{I_y}}\,|\fminus-\fplus|.
 \label{eq:measured-y-rdt-ratio}
\end{align}
The corresponding normalized-momentum ratios are
\begin{align}
 \frac{\mathcal A_{\widetilde p_x}(Q_2)}{\mathcal A_{\widetilde p_x}(Q_1)}
 &=2\sqrt{\frac{I_y}{I_x}}\,|\fminus+\fplus^{*}|,
 \label{eq:measured-px-rdt-ratio}\\
 \frac{\mathcal A_{\widetilde p_y}(Q_1)}{\mathcal A_{\widetilde p_y}(Q_2)}
 &=2\sqrt{\frac{I_x}{I_y}}\,|\fminus+\fplus|.
 \label{eq:measured-py-rdt-ratio}
\end{align}

\subsection{Mais--Ripken and RDT amplitude identities}

From \cref{eq:mr-line-amplitudes},
\begin{align}
 \frac{\mathcal A_x(Q_2)}{\mathcal A_x(Q_1)}
 &=\sqrt{\frac{J_2}{J_1}}
   \sqrt{\frac{\beta_{x2}}{\beta_{x1}}},\\
 \frac{\mathcal A_y(Q_1)}{\mathcal A_y(Q_2)}
 &=\sqrt{\frac{J_1}{J_2}}
   \sqrt{\frac{\beta_{y1}}{\beta_{y2}}}.
\end{align}
Identifying $J_1=I_x$ and $J_2=I_y$ to the working order gives
\begin{align}
 \boxed{
 r_x\equiv\frac12\sqrt{\frac{\beta_{x2}}{\beta_{x1}}}
 =|f_{1001}-f_{1010}^{*}|
 }
 \label{eq:rx-main}\\
 \boxed{
 r_y\equiv\frac12\sqrt{\frac{\beta_{y1}}{\beta_{y2}}}
 =|f_{1001}-f_{1010}|
 .}
 \label{eq:ry-main}
\end{align}
These are exact identities for the matrix-equivalent RDTs defined in \cref{sec:matrix-rdt}; for the source-summed perturbative RDTs they hold to first order in the coupling.

The normalized momentum in plane $q$ is
\begin{equation}
 \widetilde p_q=\frac{\alpha_q q+\beta_q p_q}{\sqrt{\beta_q}}.
\end{equation}
Using the Mais--Ripken projection matrices, define
\begin{align}
 r_{\widetilde p_x}
 &\equiv\frac12
 \sqrt{
  \frac{
   \alpha_x^2\beta_{x2}
   -2\alpha_x\beta_x\alpha_{x2}
   +\beta_x^2\gamma_{x2}
  }{\beta_{x1}}
 },
 \label{eq:rpx-mr}\\
 r_{\widetilde p_y}
 &\equiv\frac12
 \sqrt{
  \frac{
   \alpha_y^2\beta_{y1}
   -2\alpha_y\beta_y\alpha_{y1}
   +\beta_y^2\gamma_{y1}
  }{\beta_{y2}}
 }.
 \label{eq:rpy-mr}
\end{align}
Then
\begin{align}
 \boxed{r_{\widetilde p_x}=|f_{1001}+f_{1010}^{*}|,}
 \label{eq:rpx-rdt}\\
 \boxed{r_{\widetilde p_y}=|f_{1001}+f_{1010}|.}
 \label{eq:rpy-rdt}
\end{align}
Thus the position projection selects the difference of the two complex coupling terms, while the normalized-momentum projection selects their sum.

The interference structure is explicit:
\begin{align}
 r_x^2&=|\fminus|^2+|\fplus|^2-2\Ree(\fminus\fplus),\\
 r_{\widetilde p_x}^2&=|\fminus|^2+|\fplus|^2+2\Ree(\fminus\fplus),\\
 r_y^2&=|\fminus|^2+|\fplus|^2-2\Ree(\fminus\fplus^{*}),\\
 r_{\widetilde p_y}^2&=|\fminus|^2+|\fplus|^2+2\Ree(\fminus\fplus^{*}).
 \label{eq:interference-identities}
\end{align}
Consequently, even a relatively small $f_{1010}$ produces a first-order modulation of a Mais--Ripken position ratio around $|f_{1001}|$.

\subsection{Physical canonical-momentum ratios}

The Mais--Ripken functions give directly
\begin{align}
 \frac{\mathcal A_{p_x}(Q_2)}{\mathcal A_{p_x}(Q_1)}
 &=\sqrt{\frac{J_2}{J_1}}
   \sqrt{\frac{\gamma_{x2}}{\gamma_{x1}}},\\
 \frac{\mathcal A_{p_y}(Q_1)}{\mathcal A_{p_y}(Q_2)}
 &=\sqrt{\frac{J_1}{J_2}}
   \sqrt{\frac{\gamma_{y1}}{\gamma_{y2}}}.
\end{align}
Because $p_q=(\widetilde p_q-\alpha_q\widetilde q)/\sqrt{\beta_q}$, the pure RDT combinations are mixed by $\alpha_q$.  The corresponding lattice ratios are
\begin{align}
 \boxed{
 \frac12\sqrt{\frac{\gamma_{x2}}{\gamma_{x1}}}
 =\frac{
  \left|\fminus+\fplus^{*}+\ii\alpha_x(\fminus-\fplus^{*})\right|
 }{\sqrt{1+\alpha_x^2}}
 }
 \label{eq:physical-px-ratio}\\
 \boxed{
 \frac12\sqrt{\frac{\gamma_{y1}}{\gamma_{y2}}}
 =\frac{
  \left|\fminus+\fplus-\ii\alpha_y(\fminus-\fplus)\right|
 }{\sqrt{1+\alpha_y^2}}
 .}
 \label{eq:physical-py-ratio}
\end{align}
At an ET waist, where $\alpha_q=0$, the physical and normalized momentum ratios coincide.

\section{Exact matrix-equivalent RDTs from $W$}
\label{sec:matrix-rdt}

\subsection{Definition from the normalized coupling matrix}

The Edwards--Teng construction in \cref{sec:et-twiss} provides $g$ and the
normalized coupling matrix $\mathcal C$ directly from $W$, before any RDT is
introduced.  In the Xsuite RDT convention, the corresponding exact
matrix-equivalent linear coupling RDTs are
\begin{align}
 \boxed{
 f_{1001}=\frac{
  c_{21}-c_{12}+\ii(c_{11}+c_{22})
 }{4g}
 }
 \label{eq:f1001-from-c}\\
 \boxed{
 f_{1010}=\frac{
  c_{12}+c_{21}+\ii(c_{11}-c_{22})
 }{4g}.
 }
 \label{eq:f1010-from-c}
\end{align}
Relative to the convention used in the original matrix--Hamiltonian formulas of Calaga, Tom\'as, and Franchi~\cite{CalagaTomasFranchi}, the Xsuite convention applies an overall minus complex conjugation to each of these two RDTs.

Conversely,
\begin{equation}
 \boxed{
 \mathcal C=2g
 \begin{pmatrix}
  \Imm(\fminus+\fplus)&\Ree(\fplus-\fminus)\\
  \Ree(\fminus+\fplus)&\Imm(\fminus-\fplus)
 \end{pmatrix}.
 }
 \label{eq:c-from-rdts}
\end{equation}
Its complex row combinations are
\begin{align}
 c_{11}+\ii c_{12}&=-2\ii g(\fminus-\fplus^{*}),
 \label{eq:c-row1-complex}\\
 c_{21}+\ii c_{22}&= 2g(\fminus+\fplus^{*}),
 \label{eq:c-row2-complex}\\
 (-\bar{\mathcal C})_{11}
 +\ii(-\bar{\mathcal C})_{12}
 &=-2\ii g(\fminus^{*}-\fplus^{*}),
 \label{eq:cbar-row1-complex}\\
 (-\bar{\mathcal C})_{21}
 +\ii(-\bar{\mathcal C})_{22}
 &=2g(\fminus^{*}+\fplus^{*}).
 \label{eq:cbar-row2-complex}
\end{align}
These four identities are the matrix counterpart of \cref{eq:xtilde-rdt,eq:pxtilde-rdt,eq:ytilde-rdt,eq:pytilde-rdt}.

Equation~\eqref{eq:g-c-condition} becomes
\begin{equation}
 \det\mathcal C=4g^2\left(|\fminus|^2-|\fplus|^2\right),
 \qquad
 \boxed{
 g^2=\frac{1}{1+4(|\fminus|^2-|\fplus|^2)}.
 }
 \label{eq:g-from-rdts}
\end{equation}
Thus $g=1+\mathcal O(\fminus^2,\fplus^2)$, explaining why it is absent in the
first-order orbit formulas while canceling exactly from the amplitude ratios.
Combining \cref{eq:sigmax-factorization,eq:sigmay-factorization} with
\cref{eq:c-row1-complex,eq:c-row2-complex,eq:cbar-row1-complex,eq:cbar-row2-complex}
immediately yields the exact amplitude identities in
\cref{eq:rx-main,eq:ry-main,eq:rpx-rdt,eq:rpy-rdt}.

\subsection{Direct extraction of the coupling RDTs from the normal-form matrix}
\label{sec:direct-rdts-from-w}

The matrix-equivalent RDTs can be obtained directly from $W$ without first
reconstructing all four real entries of $\mathcal C$.  Assume the coordinate
order $(x,p_x,y,p_y)$ and that columns 3 and 4 of $W$ belong to the
vertical-like mode 2.  As stated in \cref{eq:et-rephasing-of-w}, $W$ is already
assumed to be in the ET phase gauge.

The complex mode-2 eigenvector is therefore
\begin{equation}
 \mat v_2=W_{:,3}+\ii W_{:,4},
 \qquad v_{2,y}>0,
 \label{eq:v2-et-phase-gauge}
\end{equation}
with a real, positive physical $y$ component.  The mode-1 phase does not enter
the following calculation because its Mais--Ripken functions are phase invariant.

Next obtain $g$ from the principal horizontal Mais--Ripken functions and recover the
horizontal ET Twiss parameters:
\begin{align}
 g^2
 &=\sqrt{\beta_{x1}\gamma_{x1}-\alpha_{x1}^2},
 &g&=\sqrt{g^2},
 \label{eq:direct-g-from-w}\\
 \beta_x&=\frac{\beta_{x1}}{g^2},
 &\alpha_x&=\frac{\alpha_{x1}}{g^2}.
 \label{eq:direct-et-twiss-from-w}
\end{align}
The quantities $\beta_{x1}$, $\alpha_{x1}$, and $\gamma_{x1}$ are obtained
from the first two columns of $W$ through
\cref{eq:mr-functions-from-w}.

Apply the inverse horizontal ET normalizer to the horizontal components of
$\mat v_2$:
\begin{equation}
 \begin{pmatrix}u_x\\u_{p_x}\end{pmatrix}
 =A_x^{-1}
 \begin{pmatrix}v_{2,x}\\v_{2,p_x}\end{pmatrix}
 =\begin{pmatrix}
  v_{2,x}/\sqrt{\beta_x}\\[0.3em]
  (\alpha_xv_{2,x}+\beta_xv_{2,p_x})/\sqrt{\beta_x}
 \end{pmatrix}.
 \label{eq:direct-normalized-v2}
\end{equation}
By \cref{eq:w-blocks-normalized}, $A_x^{-1}W_{x2}=\mathcal C$, so these two
complex numbers are
\begin{equation}
 u_x=c_{11}+\ii c_{12},
 \qquad
 u_{p_x}=c_{21}+\ii c_{22}.
 \label{eq:direct-c-rows}
\end{equation}
Substitution into \cref{eq:f1001-from-c,eq:f1010-from-c} finally gives
\begin{align}
 \boxed{
 f_{1001}=\frac{u_{p_x}+\ii u_x}{4g}
 }
 \label{eq:direct-f1001-from-w}\\
 \boxed{
 f_{1010}=\frac{(u_{p_x}-\ii u_x)^{*}}{4g}.
 }
 \label{eq:direct-f1010-from-w}
\end{align}
This is the direct route used by an implementation based on the complex
mode-2 eigenvector of $W$.
